\chapter{Literature Review}
\label{ch:lit_review}

The application of deep reinforcement learning to algorithmic trading and portfolio management has developed rapidly across policy optimisation, sequence modeling, risk-aware rewards, and evaluation standards. This chapter presents a focused review of the core literature directly motivating our research.

\noindent\textbf{Schulman et al. (2017)} \citep{schulman2017proximal} developed Proximal Policy Optimisation (PPO), introducing a clipped surrogate objective that restricts policy update ratios within an $\epsilon$-neighbourhood around the previous policy. This provides stable gradient updates in continuous control environments without the excessive hyperparameter sensitivity of earlier policy gradient methods. While PPO provides an effective actor--critic foundation, its application to continuous multi-asset portfolio management requires adapting the action distribution to the portfolio simplex and formulating risk-sensitive rewards.

\noindent\textbf{Zou et al. (2024)} \citep{zou2024novel} proposed an equity trading architecture that combines a recurrent LSTM encoder with PPO policy optimization. Using a 30-day historical window and a 512-dimensional hidden representation, their system showed that multi-day sequential features capture price momentum more effectively than single-day snapshots. However, their policy rewards focused on nominal trading profit for individual equities rather than risk-adjusted continuous portfolio allocation across a multi-asset universe.

\noindent\textbf{Millea (2021)} \citep{millea2021deep} provided a survey of reinforcement learning in quantitative trading, emphasizing the need for standardised evaluation protocols and risk-sensitive objective functions. The author reviewed the Differential Sharpe Ratio (DSR), which updates running first and second return moments recursively to reward risk-adjusted returns online. While DSR provides an analytical alternative to raw returns, its empirical interaction with recurrent encoders and execution frictions under standard actor--critic methods warrants systematic evaluation.

\noindent\textbf{Liu et al. (2024)} \citep{liu2024dynamic} introduced the FinRL-Meta framework, establishing dynamic market environments and a rolling walk-forward evaluation protocol. Their rolling train--validation--test scheme advances windows chronologically across market cycles, directly mitigating look-ahead bias and macroeconomic regime shifts. Although their work established rigorous financial testing standards, their continuous portfolio experiments evaluated fixed configurations without isolating the separate impacts of sequence encoders or turnover penalties.

\noindent\textbf{Wang and Liu (2025)} \citep{wang2025risk} developed a risk-sensitive reinforcement learning framework for portfolio allocation across volatile market regimes, incorporating downside risk penalties and turnover control into policy learning. Their results confirmed that penalising volatility and trading frictions preserves capital during market downturns. However, because their experiments modified the sequence backbone, reward function, and trading constraints simultaneously, the individual contribution of each mechanism was not isolated.

\noindent\textbf{Yang et al. (2022)} \citep{yang2022selective} introduced a continuous portfolio allocation algorithm parameterised over a multivariate Dirichlet distribution. Unlike heuristic Gaussian clipping or post-hoc Softmax normalization, sampling actions from a Dirichlet distribution inherently constrains target portfolio allocations to the unit simplex manifold. While their study demonstrated the mathematical appeal of Dirichlet policies in off-policy trading, it did not investigate recurrent temporal representations or explicit transaction-cost regularization.
